\documentclass[11pt]{article}
\usepackage{amsmath,amssymb,amsthm}
\usepackage[margin=1.1in]{geometry}

\newtheorem{remark}{Remark}

\title{Lagrange Multipliers, Walked Through Once:\\
How the Constraint Produces the $/n$ in the Multinoulli MLE}
\author{Yuval Lavie}
\date{\today}

\begin{document}
\maketitle

\noindent
A companion to the optimization note's Section 3: the same theory, but
executed slowly on one example --- the multinoulli MLE --- with the
role of each equation made explicit. Read this when you remember
\emph{that} the Lagrangian works but not \emph{why} the constraint
enters at the end.

\section{The problem, and why plain stationarity fails}

Maximize the multinoulli log-likelihood over the probability vector:
\begin{equation}
  \max_{p}\;\; \ell(p) = \sum_{j=1}^{k} n_j \ln p_j
  \qquad\text{subject to}\qquad
  g(p) = \sum_{j=1}^{k} p_j - 1 = 0 .
  \label{eq:problem}
\end{equation}
Try the unconstrained move first:
$\partial \ell / \partial p_j = n_j / p_j > 0$ everywhere. The gradient
is \emph{never} zero --- $\ell$ simply wants every $p_j$ larger, without
bound. There is no interior stationary point; only the constraint stops
the climb. The maximum therefore sits \emph{on} the constraint surface,
at a point where $\nabla\ell \ne 0$. Fermat's rule does not apply and
something must replace it.

\section{The idea: parallel gradients}

At a constrained maximum, no small move \emph{along the constraint
surface} can improve $\ell$ (otherwise we would make it). So the
component of $\nabla\ell$ tangent to the surface is zero ---
$\nabla\ell$ is \emph{perpendicular} to the surface. But $\nabla g$ is
also perpendicular to the surface: the gradient of a function is always
normal to its level sets, and the surface \emph{is} the level set
$g = 0$. Two vectors perpendicular to the same surface are parallel:
\begin{equation}
  \boxed{\;\nabla \ell(p^\ast) \;=\; \lambda\, \nabla g(p^\ast)\;}
  \label{eq:parallel}
\end{equation}
for some scalar $\lambda$ --- the \emph{Lagrange multiplier}. That one
line is the entire theory. The \emph{Lagrangian}
\begin{equation}
  \Lambda(p, \lambda)
  \;=\; \ell(p) \;+\; \lambda\Bigl(1 - \sum_j p_j\Bigr)
  \label{eq:lagrangian}
\end{equation}
is bookkeeping that packages \eqref{eq:parallel}: demanding
$\partial \Lambda / \partial p_j = 0$ for all $j$ reproduces the
parallel-gradients condition, and demanding
$\partial \Lambda / \partial \lambda = 0$ reproduces the constraint
itself. One unconstrained stationarity problem in $(p, \lambda)$
encodes the constrained problem in $p$.

\section{The mechanics, step by step}

\paragraph{Step 1: stationarity in each $p_j$ --- the shape.}
\begin{equation*}
  \frac{\partial \Lambda}{\partial p_j}
  = \frac{n_j}{p_j} - \lambda = 0
  \qquad\Longrightarrow\qquad
  p_j = \frac{n_j}{\lambda}.
\end{equation*}
Notice what this delivers: the \emph{shape} of the solution --- each
$p_j$ proportional to its count --- with one unknown scale $\lambda$
left dangling. This is typical. The gradient equations determine the
answer only up to exactly the freedom the constraint will remove.

\paragraph{Step 2: the constraint fixes the scale --- the $/n$.}
We hold $k$ expressions $p_j = n_j/\lambda$ with one unknown, and one
equation not yet used: $\sum_j p_j = 1$. Substitute:
\begin{equation*}
  \sum_{j} \frac{n_j}{\lambda} = 1
  \qquad\Longrightarrow\qquad
  \frac{1}{\lambda}\underbrace{\sum_{j} n_j}_{=\,n} = 1
  \qquad\Longrightarrow\qquad
  \lambda = n .
\end{equation*}
\emph{This} is where the constraint acts: not in the differentiation,
but as the closing equation that pins the multiplier.

\paragraph{Step 3: back-substitute.}
\begin{equation}
  \hat p_j \;=\; \frac{n_j}{n}\,.
  \label{eq:answer}
\end{equation}

\medskip
\noindent
The division of labor, in one sentence: \textbf{the gradient equations
pin the direction of the answer; the constraint pins its scale.}

\section{What $\lambda$ means, and why we discard it}

The multiplier is the \emph{exchange rate} between objective and
constraint: at the optimum,
\begin{equation}
  \lambda \;=\; \frac{d\,\ell^\ast}{d\,c}
  \qquad\text{for the relaxed constraint } \textstyle\sum_j p_j = c ,
  \label{eq:shadow}
\end{equation}
the improvement in optimal log-likelihood per unit of loosened
constraint (``shadow price''). Here $\lambda = n$: allowing the
probabilities to sum to $1+\varepsilon$ would buy roughly
$n\varepsilon$ extra log-likelihood --- each of the $n$ observations
would enjoy proportionally inflated probability. We solve for
$\lambda$, use it, and throw it away: it is scaffolding, not part of
the answer.

\begin{remark}[The certificate still matters]
Parallel gradients \eqref{eq:parallel} is only \emph{necessary} --- it
holds equally at constrained minima and saddles. Here
$\ell(p) = \sum_j n_j \ln p_j$ is strictly concave and the simplex is
convex, so the unique point satisfying \eqref{eq:parallel} is the
global constrained maximum. Without concavity, check candidates.
\end{remark}

\begin{remark}[The constraint we ignored]
Problem \eqref{eq:problem} also carries $p_j \ge 0$, silently. We could
ignore it because $\ln p_j \to -\infty$ as $p_j \to 0$ whenever
$n_j > 0$: the objective polices the boundary by itself. When
inequality constraints \emph{can} bind, equality Lagrange theory
generalizes to the KKT conditions --- multipliers that must be
nonnegative and vanish on inactive constraints --- covered in the
optimization note, Section 3.
\end{remark}

\end{document}
